% Generated by scripts/spec_latex.py; edit the TLA+ source.
\batchmode %% Suppresses most terminal output.
\documentclass{article}
\setlength{\textwidth}{500pt}
\setlength{\textheight}{680pt}
\addtolength{\hoffset}{-70pt}
\addtolength{\voffset}{-60pt}
\usepackage{latexsym}
\usepackage{ifthen}
% \usepackage{color}
%%%%%%%%%%%%%%%%%%%%%%%%%%%%%%%%%%%%%%%%%%%%%%%%%%%%%%%%%%%%%%%%%%%%%%%%%%%%%
% SWITCHES                                                                  %
%%%%%%%%%%%%%%%%%%%%%%%%%%%%%%%%%%%%%%%%%%%%%%%%%%%%%%%%%%%%%%%%%%%%%%%%%%%%%
\newboolean{shading} 
\setboolean{shading}{false}
\makeatletter
 %% this is needed only when inserted into the file, not when
 %% used as a package file.
%%%%%%%%%%%%%%%%%%%%%%%%%%%%%%%%%%%%%%%%%%%%%%%%%%%%%%%%%%%%%%%%%%%%%%%%%%%%%
%                                                                           %
% DEFINITIONS OF SYMBOL-PRODUCING COMMANDS                                  %
%                                                                           %
%    TLA+      LaTeX                                                        %
%    symbol    command                                                      %
%    ------    -------                                                      %
%    =>        \implies                                                     %
%    <:        \ltcolon                                                     %
%    :>        \colongt                                                     %
%    ==        \defeq                                                       %
%    ..        \dotdot                                                      %
%    ::        \coloncolon                                                  %
%    =|        \eqdash                                                      %
%    ++        \pp                                                          %
%    --        \mm                                                          %
%    **        \stst                                                        %
%    //        \slsl                                                        %
%    ^         \ct                                                          %
%    \A        \A                                                           %
%    \E        \E                                                           %
%    \AA       \AA                                                          %
%    \EE       \EE                                                          %
%%%%%%%%%%%%%%%%%%%%%%%%%%%%%%%%%%%%%%%%%%%%%%%%%%%%%%%%%%%%%%%%%%%%%%%%%%%%%
\newlength{\symlength}
\newcommand{\implies}{\Rightarrow}
\newcommand{\ltcolon}{\mathrel{<\!\!\mbox{:}}}
\newcommand{\colongt}{\mathrel{\!\mbox{:}\!\!>}}
\newcommand{\defeq}{\;\mathrel{\smash   %% keep this symbol from being too tall
    {{\stackrel{\scriptscriptstyle\Delta}{=}}}}\;}
\newcommand{\dotdot}{\mathrel{\ldotp\ldotp}}
\newcommand{\coloncolon}{\mathrel{::\;}}
\newcommand{\eqdash}{\mathrel = \joinrel \hspace{-.28em}|}
\newcommand{\pp}{\mathbin{++}}
\newcommand{\mm}{\mathbin{--}}
\newcommand{\stst}{*\!*}
\newcommand{\slsl}{/\!/}
\newcommand{\ct}{\hat{\hspace{.4em}}}
\newcommand{\A}{\forall}
\newcommand{\E}{\exists}
\renewcommand{\AA}{\makebox{$\raisebox{.05em}{\makebox[0pt][l]{%
   $\forall\hspace{-.517em}\forall\hspace{-.517em}\forall$}}%
   \forall\hspace{-.517em}\forall \hspace{-.517em}\forall\,$}}
\newcommand{\EE}{\makebox{$\raisebox{.05em}{\makebox[0pt][l]{%
   $\exists\hspace{-.517em}\exists\hspace{-.517em}\exists$}}%
   \exists\hspace{-.517em}\exists\hspace{-.517em}\exists\,$}}
\newcommand{\whileop}{\.{\stackrel
  {\mbox{\raisebox{-.3em}[0pt][0pt]{$\scriptscriptstyle+\;\,$}}}%
  {-\hspace{-.16em}\triangleright}}}

% Commands are defined to produce the upper-case keywords.
% Note that some have space after them.
\newcommand{\ASSUME}{\textsc{assume }}
\newcommand{\ASSUMPTION}{\textsc{assumption }}
\newcommand{\AXIOM}{\textsc{axiom }}
\newcommand{\BOOLEAN}{\textsc{boolean }}
\newcommand{\CASE}{\textsc{case }}
\newcommand{\CONSTANT}{\textsc{constant }}
\newcommand{\CONSTANTS}{\textsc{constants }}
\newcommand{\ELSE}{\settowidth{\symlength}{\THEN}%
   \makebox[\symlength][l]{\textsc{ else}}}
\newcommand{\EXCEPT}{\textsc{ except }}
\newcommand{\EXTENDS}{\textsc{extends }}
\newcommand{\FALSE}{\textsc{false}}
\newcommand{\IF}{\textsc{if }}
\newcommand{\IN}{\settowidth{\symlength}{\LET}%
   \makebox[\symlength][l]{\textsc{in}}}
\newcommand{\INSTANCE}{\textsc{instance }}
\newcommand{\LET}{\textsc{let }}
\newcommand{\LOCAL}{\textsc{local }}
\newcommand{\MODULE}{\textsc{module }}
\newcommand{\OTHER}{\textsc{other }}
\newcommand{\STRING}{\textsc{string}}
\newcommand{\THEN}{\textsc{ then }}
\newcommand{\THEOREM}{\textsc{theorem }}
\newcommand{\LEMMA}{\textsc{lemma }}
\newcommand{\PROPOSITION}{\textsc{proposition }}
\newcommand{\COROLLARY}{\textsc{corollary }}
\newcommand{\TRUE}{\textsc{true}}
\newcommand{\VARIABLE}{\textsc{variable }}
\newcommand{\VARIABLES}{\textsc{variables }}
\newcommand{\WITH}{\textsc{ with }}
\newcommand{\WF}{\textrm{WF}}
\newcommand{\SF}{\textrm{SF}}
\newcommand{\CHOOSE}{\textsc{choose }}
\newcommand{\ENABLED}{\textsc{enabled }}
\newcommand{\UNCHANGED}{\textsc{unchanged }}
\newcommand{\SUBSET}{\textsc{subset }}
\newcommand{\UNION}{\textsc{union }}
\newcommand{\DOMAIN}{\textsc{domain }}
% Added for tla2tex
\newcommand{\BY}{\textsc{by }}
\newcommand{\OBVIOUS}{\textsc{obvious }}
\newcommand{\HAVE}{\textsc{have }}
\newcommand{\QED}{\textsc{qed }}
\newcommand{\TAKE}{\textsc{take }}
\newcommand{\DEF}{\textsc{ def }}
\newcommand{\HIDE}{\textsc{hide }}
\newcommand{\RECURSIVE}{\textsc{recursive }}
\newcommand{\USE}{\textsc{use }}
\newcommand{\DEFINE}{\textsc{define }}
\newcommand{\PROOF}{\textsc{proof }}
\newcommand{\WITNESS}{\textsc{witness }}
\newcommand{\PICK}{\textsc{pick }}
\newcommand{\DEFS}{\textsc{defs }}
\newcommand{\PROVE}{\settowidth{\symlength}{\ASSUME}%
   \makebox[\symlength][l]{\textsc{prove}}\@s{-4.1}}%
  %% The \@s{-4.1) is a kludge added on 24 Oct 2009 [happy birthday, Ellen]
  %% so the correct alignment occurs if the user types
  %%   ASSUME X
  %%   PROVE  Y
  %% because it cancels the extra 4.1 pts added because of the 
  %% extra space after the PROVE.  This seems to works OK.
  %% However, the 4.1 equals Parameters.LaTeXLeftSpace(1) and
  %% should be changed if that method ever changes.
\newcommand{\SUFFICES}{\textsc{suffices }}
\newcommand{\NEW}{\textsc{new }}
\newcommand{\LAMBDA}{\textsc{lambda }}
\newcommand{\STATE}{\textsc{state }}
\newcommand{\ACTION}{\textsc{action }}
\newcommand{\TEMPORAL}{\textsc{temporal }}
\newcommand{\ONLY}{\textsc{only }}              %% added by LL on 2 Oct 2009
\newcommand{\OMITTED}{\textsc{omitted }}        %% added by LL on 31 Oct 2009
\newcommand{\@pfstepnum}[2]{\ensuremath{\langle#1\rangle}\textrm{#2}}
\newcommand{\bang}{\@s{1}\mbox{\small !}\@s{1}}
%% We should format || differently in PlusCal code than in TLA+ formulas.
\newcommand{\p@barbar}{\ifpcalsymbols
   \,\,\rule[-.25em]{.075em}{1em}\hspace*{.2em}\rule[-.25em]{.075em}{1em}\,\,%
   \else \,||\,\fi}
%% PlusCal keywords
\newcommand{\p@fair}{\textbf{fair }}
\newcommand{\p@semicolon}{\textbf{\,; }}
\newcommand{\p@algorithm}{\textbf{algorithm }}
\newcommand{\p@mmfair}{\textbf{-{}-fair }}
\newcommand{\p@mmalgorithm}{\textbf{-{}-algorithm }}
\newcommand{\p@assert}{\textbf{assert }}
\newcommand{\p@await}{\textbf{await }}
\newcommand{\p@begin}{\textbf{begin }}
\newcommand{\p@end}{\textbf{end }}
\newcommand{\p@call}{\textbf{call }}
\newcommand{\p@define}{\textbf{define }}
\newcommand{\p@do}{\textbf{ do }}
\newcommand{\p@either}{\textbf{either }}
\newcommand{\p@or}{\textbf{or }}
\newcommand{\p@goto}{\textbf{goto }}
\newcommand{\p@if}{\textbf{if }}
\newcommand{\p@then}{\,\,\textbf{then }}
\newcommand{\p@else}{\ifcsyntax \textbf{else } \else \,\,\textbf{else }\fi}
\newcommand{\p@elsif}{\,\,\textbf{elsif }}
\newcommand{\p@macro}{\textbf{macro }}
\newcommand{\p@print}{\textbf{print }}
\newcommand{\p@procedure}{\textbf{procedure }}
\newcommand{\p@process}{\textbf{process }}
\newcommand{\p@return}{\textbf{return}}
\newcommand{\p@skip}{\textbf{skip}}
\newcommand{\p@variable}{\textbf{variable }}
\newcommand{\p@variables}{\textbf{variables }}
\newcommand{\p@while}{\textbf{while }}
\newcommand{\p@when}{\textbf{when }}
\newcommand{\p@with}{\textbf{with }}
\newcommand{\p@lparen}{\textbf{(\,\,}}
\newcommand{\p@rparen}{\textbf{\,\,) }}   
\newcommand{\p@lbrace}{\textbf{\{\,\,}}   
\newcommand{\p@rbrace}{\textbf{\,\,\} }}

%%%%%%%%%%%%%%%%%%%%%%%%%%%%%%%%%%%%%%%%%%%%%%%%%%%%%%%%%
% REDEFINE STANDARD COMMANDS TO MAKE THEM FORMAT BETTER %
%                                                       %
% We redefine \in and \notin                            %
%%%%%%%%%%%%%%%%%%%%%%%%%%%%%%%%%%%%%%%%%%%%%%%%%%%%%%%%%
\renewcommand{\_}{\rule{.4em}{.06em}\hspace{.05em}}
\newlength{\equalswidth}
\let\oldin=\in
\let\oldnotin=\notin
\renewcommand{\in}{%
   {\settowidth{\equalswidth}{$\.{=}$}\makebox[\equalswidth][c]{$\oldin$}}}
\renewcommand{\notin}{%
   {\settowidth{\equalswidth}{$\.{=}$}\makebox[\equalswidth]{$\oldnotin$}}}


%%%%%%%%%%%%%%%%%%%%%%%%%%%%%%%%%%%%%%%%%%%%%%%%%%%%
%                                                  %
% HORIZONTAL BARS:                                 %
%                                                  %
%   \moduleLeftDash    |~~~~~~~~~~                 %
%   \moduleRightDash    ~~~~~~~~~~|                %
%   \midbar            |----------|                %
%   \bottombar         |__________|                %
%%%%%%%%%%%%%%%%%%%%%%%%%%%%%%%%%%%%%%%%%%%%%%%%%%%%
\newlength{\charwidth}\settowidth{\charwidth}{{\small\tt M}}
\newlength{\boxrulewd}\setlength{\boxrulewd}{.4pt}
\newlength{\boxlineht}\setlength{\boxlineht}{.5\baselineskip}
\newcommand{\boxsep}{\charwidth}
\newlength{\boxruleht}\setlength{\boxruleht}{.5ex}
\newlength{\boxruledp}\setlength{\boxruledp}{-\boxruleht}
\addtolength{\boxruledp}{\boxrulewd}
\newcommand{\boxrule}{\leaders\hrule height \boxruleht depth \boxruledp
                      \hfill\mbox{}}
\newcommand{\@computerule}{%
  \setlength{\boxruleht}{.5ex}%
  \setlength{\boxruledp}{-\boxruleht}%
  \addtolength{\boxruledp}{\boxrulewd}}

\newcommand{\bottombar}{\hspace{-\boxsep}%
  \raisebox{-\boxrulewd}[0pt][0pt]{\rule[.5ex]{\boxrulewd}{\boxlineht}}%
  \boxrule
  \raisebox{-\boxrulewd}[0pt][0pt]{%
      \rule[.5ex]{\boxrulewd}{\boxlineht}}\hspace{-\boxsep}\vspace{0pt}}

\newcommand{\moduleLeftDash}%
   {\hspace*{-\boxsep}%
     \raisebox{-\boxlineht}[0pt][0pt]{\rule[.5ex]{\boxrulewd
               }{\boxlineht}}%
    \boxrule\hspace*{.4em }}

\newcommand{\moduleRightDash}%
    {\hspace*{.4em}\boxrule
    \raisebox{-\boxlineht}[0pt][0pt]{\rule[.5ex]{\boxrulewd
               }{\boxlineht}}\hspace{-\boxsep}}%\vspace{.2em}

\newcommand{\midbar}{\hspace{-\boxsep}\raisebox{-.5\boxlineht}[0pt][0pt]{%
   \rule[.5ex]{\boxrulewd}{\boxlineht}}\boxrule\raisebox{-.5\boxlineht%
   }[0pt][0pt]{\rule[.5ex]{\boxrulewd}{\boxlineht}}\hspace{-\boxsep}}

%%%%%%%%%%%%%%%%%%%%%%%%%%%%%%%%%%%%%%%%%%%%%%%%%%%%%%%%%%%%%%%%%%%%%%%%%%%%%
% FORMATING COMMANDS                                                        %
%%%%%%%%%%%%%%%%%%%%%%%%%%%%%%%%%%%%%%%%%%%%%%%%%%%%%%%%%%%%%%%%%%%%%%%%%%%%%

%%%%%%%%%%%%%%%%%%%%%%%%%%%%%%%%%%%%%%%%%%%%%%%%%%%%%%%%%%%%%%%%%%%%%%%%%%%%%
% PLUSCAL SHADING                                                           %
%%%%%%%%%%%%%%%%%%%%%%%%%%%%%%%%%%%%%%%%%%%%%%%%%%%%%%%%%%%%%%%%%%%%%%%%%%%%%

% The TeX pcalshading switch is set on to cause PlusCal shading to be
% performed.  This changes the behavior of the following commands and
% environments to cause full-width shading to be performed on all lines.
% 
%   \tstrut \@x cpar mcom \@pvspace
% 
% The TeX pcalsymbols switch is turned on when typesetting a PlusCal algorithm,
% whether or not shading is being performed.  It causes symbols (other than
% parentheses and braces and PlusCal-only keywords) that should be typeset
% differently depending on whether they are in an algorithm to be typeset
% appropriately.  Currently, the only such symbol is "||".
%
% The TeX csyntax switch is turned on when typesetting a PlusCal algorithm in
% c-syntax.  This allows symbols to be format differently in the two syntaxes.
% The "else" keyword is the only one that is.

\newif\ifpcalshading \pcalshadingfalse
\newif\ifpcalsymbols \pcalsymbolsfalse
\newif\ifcsyntax     \csyntaxtrue

% The \@pvspace command makes a vertical space.  It uses \vspace
% except with \ifpcalshading, in which case it sets \pvcalvspace
% and the space is added by a following \@x command.
%
\newlength{\pcalvspace}\setlength{\pcalvspace}{0pt}%
\newcommand{\@pvspace}[1]{%
  \ifpcalshading
     \par\global\setlength{\pcalvspace}{#1}%
  \else
     \par\vspace{#1}%
  \fi
}

% The lcom environment was changed to set \lcomindent equal to
% the indentation it produces.  This length is used by the
% cpar environment to make shading extend for the full width
% of the line.  This assumes that lcom environments are not
% nested.  I hope TLATeX does not nest them.
%
\newlength{\lcomindent}%
\setlength{\lcomindent}{0pt}%

%\tstrut: A strut to produce inter-paragraph space in a comment.
%\rstrut: A strut to extend the bottom of a one-line comment so
%         there's no break in the shading between comments on 
%         successive lines.
\newcommand\tstrut%
  {\raisebox{\vshadelen}{\raisebox{-.25em}{\rule{0pt}{1.15em}}}%
   \global\setlength{\vshadelen}{0pt}}
\newcommand\rstrut{\raisebox{-.25em}{\rule{0pt}{1.15em}}%
 \global\setlength{\vshadelen}{0pt}}


% \.{op} formats operator op in math mode with empty boxes on either side.
% Used because TeX otherwise vary the amount of space it leaves around op.
\renewcommand{\.}[1]{\ensuremath{\mbox{}#1\mbox{}}}

% \@s{n} produces an n-point space
\newcommand{\@s}[1]{\hspace{#1pt}}           

% \@x{txt} starts a specification line in the beginning with txt
% in the final LaTeX source.
\newlength{\@xlen}
\newcommand\xtstrut%
  {\setlength{\@xlen}{1.05em}%
   \addtolength{\@xlen}{\pcalvspace}%
    \raisebox{\vshadelen}{\raisebox{-.25em}{\rule{0pt}{\@xlen}}}%
   \global\setlength{\vshadelen}{0pt}%
   \global\setlength{\pcalvspace}{0pt}}

\newcommand{\@x}[1]{\par
  \ifpcalshading
  \makebox[0pt][l]{\shadebox{\xtstrut\hspace*{\textwidth}}}%
  \fi
  \mbox{$\mbox{}#1\mbox{}$}}  

% \@xx{txt} continues a specification line with the text txt.
\newcommand{\@xx}[1]{\mbox{$\mbox{}#1\mbox{}$}}  

% \@y{cmt} produces a one-line comment.
\newcommand{\@y}[1]{\mbox{\footnotesize\hspace{.65em}%
  \ifthenelse{\boolean{shading}}{%
      \shadebox{#1\hspace{-\the\lastskip}\rstrut}}%
               {#1\hspace{-\the\lastskip}\rstrut}}}

% \@z{cmt} produces a zero-width one-line comment.
\newcommand{\@z}[1]{\makebox[0pt][l]{\footnotesize
  \ifthenelse{\boolean{shading}}{%
      \shadebox{#1\hspace{-\the\lastskip}\rstrut}}%
               {#1\hspace{-\the\lastskip}\rstrut}}}


% \@w{str} produces the TLA+ string "str".
\newcommand{\@w}[1]{\textsf{``{#1}''}}             


%%%%%%%%%%%%%%%%%%%%%%%%%%%%%%%%%%%%%%%%%%%%%%%%%%%%%%%%%%%%%%%%%%%%%%%%%%%%%
% SHADING                                                                   %
%%%%%%%%%%%%%%%%%%%%%%%%%%%%%%%%%%%%%%%%%%%%%%%%%%%%%%%%%%%%%%%%%%%%%%%%%%%%%
\def\graymargin{1}
  % The number of points of margin in the shaded box.

% \definecolor{boxshade}{gray}{.85}
% Defines the darkness of the shading: 1 = white, 0 = black
% Added by TLATeX only if needed.

% \shadebox{txt} puts txt in a shaded box.
\newlength{\templena}
\newlength{\templenb}
\newsavebox{\tempboxa}
\newcommand{\shadebox}[1]{{\setlength{\fboxsep}{\graymargin pt}%
     \savebox{\tempboxa}{#1}%
     \settoheight{\templena}{\usebox{\tempboxa}}%
     \settodepth{\templenb}{\usebox{\tempboxa}}%
     \hspace*{-\fboxsep}\raisebox{0pt}[\templena][\templenb]%
        {\colorbox{boxshade}{\usebox{\tempboxa}}}\hspace*{-\fboxsep}}}

% \vshade{n} makes an n-point inter-paragraph space, with
%  shading if the `shading' flag is true.
\newlength{\vshadelen}
\setlength{\vshadelen}{0pt}
\newcommand{\vshade}[1]{\ifthenelse{\boolean{shading}}%
   {\global\setlength{\vshadelen}{#1pt}}%
   {\vspace{#1pt}}}

\newlength{\boxwidth}
\newlength{\multicommentdepth}

%%%%%%%%%%%%%%%%%%%%%%%%%%%%%%%%%%%%%%%%%%%%%%%%%%%%%%%%%%%%%%%%%%%%%%%%%%%%%
% THE cpar ENVIRONMENT                                                      %
% ^^^^^^^^^^^^^^^^^^^^                                                      %
% The LaTeX input                                                           %
%                                                                           %
%   \begin{cpar}{pop}{nest}{isLabel}{d}{e}{arg6}                            %
%     XXXXXXXXXXXXXXX                                                       %
%     XXXXXXXXXXXXXXX                                                       %
%     XXXXXXXXXXXXXXX                                                       %
%   \end{cpar}                                                              %
%                                                                           %
% produces one of two possible results.  If isLabel is the letter "T",      %
% it produces the following, where [label] is the result of typesetting     %
% arg6 in an LR box, and d is is a number representing a distance in        %
% points.                                                                   %
%                                                                           %
%   prevailing |<-- d -->[label]<- e ->XXXXXXXXXXXXXXX                      %
%         left |                       XXXXXXXXXXXXXXX                      %
%       margin |                       XXXXXXXXXXXXXXX                      %
%                                                                           %
% If isLabel is the letter "F", then it produces                            %
%                                                                           %
%   prevailing |<-- d -->XXXXXXXXXXXXXXXXXXXXXXX                            %
%         left |         <- e ->XXXXXXXXXXXXXXXX                            %
%       margin |                XXXXXXXXXXXXXXXX                            %
%                                                                           %
% where d and e are numbers representing distances in points.               %
%                                                                           %
% The prevailing left margin is the one in effect before the most recent    %
% pop (argument 1) cpar environments with "T" as the nest argument, where   %
% pop is a number \geq 0.                                                   %
%                                                                           %
% If the nest argument is the letter "T", then the prevailing left          %
% margin is moved to the left of the second (and following) lines of        %
% X's.  Otherwise, the prevailing left margin is left unchanged.            %
%                                                                           %
% An \unnest{n} command moves the prevailing left margin to where it was    %
% before the most recent n cpar environments with "T" as the nesting        %
% argument.                                                                 %
%                                                                           %
% The environment leaves no vertical space above or below it, or between    %
% its paragraphs.  (TLATeX inserts the proper amount of vertical space.)    %
%%%%%%%%%%%%%%%%%%%%%%%%%%%%%%%%%%%%%%%%%%%%%%%%%%%%%%%%%%%%%%%%%%%%%%%%%%%%%

\newcounter{pardepth}
\setcounter{pardepth}{0}

% \setgmargin{txt} defines \gmarginN to be txt, where N is \roman{pardepth}.
% \thegmargin equals \gmarginN, where N is \roman{pardepth}.
\newcommand{\setgmargin}[1]{%
  \expandafter\xdef\csname gmargin\roman{pardepth}\endcsname{#1}}
\newcommand{\thegmargin}{\csname gmargin\roman{pardepth}\endcsname}
\newcommand{\gmargin}{0pt}

\newsavebox{\tempsbox}

\newlength{\@cparht}
\newlength{\@cpardp}
\newenvironment{cpar}[6]{%
  \addtocounter{pardepth}{-#1}%
  \ifthenelse{\boolean{shading}}{\par\begin{lrbox}{\tempsbox}%
                                 \begin{minipage}[t]{\linewidth}}{}%
  \begin{list}{}{%
     \edef\temp{\thegmargin}
     \ifthenelse{\equal{#3}{T}}%
       {\settowidth{\leftmargin}{\hspace{\temp}\footnotesize #6\hspace{#5pt}}%
        \addtolength{\leftmargin}{#4pt}}%
       {\setlength{\leftmargin}{#4pt}%
        \addtolength{\leftmargin}{#5pt}%
        \addtolength{\leftmargin}{\temp}%
        \setlength{\itemindent}{-#5pt}}%
      \ifthenelse{\equal{#2}{T}}{\addtocounter{pardepth}{1}%
                                 \setgmargin{\the\leftmargin}}{}%
      \setlength{\labelwidth}{0pt}%
      \setlength{\labelsep}{0pt}%
      \setlength{\itemindent}{-\leftmargin}%
      \setlength{\topsep}{0pt}%
      \setlength{\parsep}{0pt}%
      \setlength{\partopsep}{0pt}%
      \setlength{\parskip}{0pt}%
      \setlength{\itemsep}{0pt}
      \setlength{\itemindent}{#4pt}%
      \addtolength{\itemindent}{-\leftmargin}}%
   \ifthenelse{\equal{#3}{T}}%
      {\item[\tstrut\footnotesize \hspace{\temp}{#6}\hspace{#5pt}]
        }%
      {\item[\tstrut\hspace{\temp}]%
         }%
   \footnotesize}
 {\hspace{-\the\lastskip}\tstrut
 \end{list}%
  \ifthenelse{\boolean{shading}}%
          {\end{minipage}%
           \end{lrbox}%
           \ifpcalshading
             \setlength{\@cparht}{\ht\tempsbox}%
             \setlength{\@cpardp}{\dp\tempsbox}%
             \addtolength{\@cparht}{.15em}%
             \addtolength{\@cpardp}{.2em}%
             \addtolength{\@cparht}{\@cpardp}%
            % I don't know what's going on here.  I want to add a
            % \pcalvspace high shaded line, but I don't know how to
            % do it.  A little trial and error shows that the following
            % does a reasonable job approximating that, eliminating
            % the line if \pcalvspace is small.
            \addtolength{\@cparht}{\pcalvspace}%
             \ifdim \pcalvspace > .8em
               \addtolength{\pcalvspace}{-.2em}%
               \hspace*{-\lcomindent}%
               \shadebox{\rule{0pt}{\pcalvspace}\hspace*{\textwidth}}\par
               \global\setlength{\pcalvspace}{0pt}%
               \fi
             \hspace*{-\lcomindent}%
             \makebox[0pt][l]{\raisebox{-\@cpardp}[0pt][0pt]{%
                 \shadebox{\rule{0pt}{\@cparht}\hspace*{\textwidth}}}}%
             \hspace*{\lcomindent}\usebox{\tempsbox}%
             \par
           \else
             \shadebox{\usebox{\tempsbox}}\par
           \fi}%
           {}%
  }

%%%%%%%%%%%%%%%%%%%%%%%%%%%%%%%%%%%%%%%%%%%%%%%%%%%%%%%%%%%%%%%%%%%%%%%%%%%%%%
% THE ppar ENVIRONMENT                                                       %
% ^^^^^^^^^^^^^^^^^^^^                                                       %
% The environment                                                            %
%                                                                            %
%   \begin{ppar} ... \end{ppar}                                              %
%                                                                            %
% is equivalent to                                                           %
%                                                                            %
%   \begin{cpar}{0}{F}{F}{0}{0}{} ... \end{cpar}                             %
%                                                                            %
% The environment is put around each line of the output for a PlusCal        %
% algorithm.                                                                 %
%%%%%%%%%%%%%%%%%%%%%%%%%%%%%%%%%%%%%%%%%%%%%%%%%%%%%%%%%%%%%%%%%%%%%%%%%%%%%%
%\newenvironment{ppar}{%
%  \ifthenelse{\boolean{shading}}{\par\begin{lrbox}{\tempsbox}%
%                                 \begin{minipage}[t]{\linewidth}}{}%
%  \begin{list}{}{%
%     \edef\temp{\thegmargin}
%        \setlength{\leftmargin}{0pt}%
%        \addtolength{\leftmargin}{\temp}%
%        \setlength{\itemindent}{0pt}%
%      \setlength{\labelwidth}{0pt}%
%      \setlength{\labelsep}{0pt}%
%      \setlength{\itemindent}{-\leftmargin}%
%      \setlength{\topsep}{0pt}%
%      \setlength{\parsep}{0pt}%
%      \setlength{\partopsep}{0pt}%
%      \setlength{\parskip}{0pt}%
%      \setlength{\itemsep}{0pt}
%      \setlength{\itemindent}{0pt}%
%      \addtolength{\itemindent}{-\leftmargin}}%
%      \item[\tstrut\hspace{\temp}]}%
% {\hspace{-\the\lastskip}\tstrut
% \end{list}%
%  \ifthenelse{\boolean{shading}}{\end{minipage}  
%                                 \end{lrbox}%
%                                 \shadebox{\usebox{\tempsbox}}\par}{}%
%  }

 %%% TESTING
 \newcommand{\xtest}[1]{\par
 \makebox[0pt][l]{\shadebox{\xtstrut\hspace*{\textwidth}}}%
 \mbox{$\mbox{}#1\mbox{}$}} 

% \newcommand{\xxtest}[1]{\par
% \makebox[0pt][l]{\shadebox{\xtstrut{#1}\hspace*{\textwidth}}}%
% \mbox{$\mbox{}#1\mbox{}$}} 

%\newlength{\pcalvspace}
%\setlength{\pcalvspace}{0pt}
% \newlength{\xxtestlen}
% \setlength{\xxtestlen}{0pt}
% \newcommand\xtstrut%
%   {\setlength{\xxtestlen}{1.15em}%
%    \addtolength{\xxtestlen}{\pcalvspace}%
%     \raisebox{\vshadelen}{\raisebox{-.25em}{\rule{0pt}{\xxtestlen}}}%
%    \global\setlength{\vshadelen}{0pt}%
%    \global\setlength{\pcalvspace}{0pt}}
   
   %%%% TESTING
   
   %% The xcpar environment
   %%  Note: overloaded use of \pcalvspace for testing.
   %%
%   \newlength{\xcparht}%
%   \newlength{\xcpardp}%
   
%   \newenvironment{xcpar}[6]{%
%  \addtocounter{pardepth}{-#1}%
%  \ifthenelse{\boolean{shading}}{\par\begin{lrbox}{\tempsbox}%
%                                 \begin{minipage}[t]{\linewidth}}{}%
%  \begin{list}{}{%
%     \edef\temp{\thegmargin}%
%     \ifthenelse{\equal{#3}{T}}%
%       {\settowidth{\leftmargin}{\hspace{\temp}\footnotesize #6\hspace{#5pt}}%
%        \addtolength{\leftmargin}{#4pt}}%
%       {\setlength{\leftmargin}{#4pt}%
%        \addtolength{\leftmargin}{#5pt}%
%        \addtolength{\leftmargin}{\temp}%
%        \setlength{\itemindent}{-#5pt}}%
%      \ifthenelse{\equal{#2}{T}}{\addtocounter{pardepth}{1}%
%                                 \setgmargin{\the\leftmargin}}{}%
%      \setlength{\labelwidth}{0pt}%
%      \setlength{\labelsep}{0pt}%
%      \setlength{\itemindent}{-\leftmargin}%
%      \setlength{\topsep}{0pt}%
%      \setlength{\parsep}{0pt}%
%      \setlength{\partopsep}{0pt}%
%      \setlength{\parskip}{0pt}%
%      \setlength{\itemsep}{0pt}%
%      \setlength{\itemindent}{#4pt}%
%      \addtolength{\itemindent}{-\leftmargin}}%
%   \ifthenelse{\equal{#3}{T}}%
%      {\item[\xtstrut\footnotesize \hspace{\temp}{#6}\hspace{#5pt}]%
%        }%
%      {\item[\xtstrut\hspace{\temp}]%
%         }%
%   \footnotesize}
% {\hspace{-\the\lastskip}\tstrut
% \end{list}%
%  \ifthenelse{\boolean{shading}}{\end{minipage}  
%                                 \end{lrbox}%
%                                 \setlength{\xcparht}{\ht\tempsbox}%
%                                 \setlength{\xcpardp}{\dp\tempsbox}%
%                                 \addtolength{\xcparht}{.15em}%
%                                 \addtolength{\xcpardp}{.2em}%
%                                 \addtolength{\xcparht}{\xcpardp}%
%                                 \hspace*{-\lcomindent}%
%                                 \makebox[0pt][l]{\raisebox{-\xcpardp}[0pt][0pt]{%
%                                      \shadebox{\rule{0pt}{\xcparht}\hspace*{\textwidth}}}}%
%                                 \hspace*{\lcomindent}\usebox{\tempsbox}%
%                                 \par}{}%
%  }
%  
% \newlength{\xmcomlen}
%\newenvironment{xmcom}[1]{%
%  \setcounter{pardepth}{0}%
%  \hspace{.65em}%
%  \begin{lrbox}{\alignbox}\sloppypar%
%      \setboolean{shading}{false}%
%      \setlength{\boxwidth}{#1pt}%
%      \addtolength{\boxwidth}{-.65em}%
%      \begin{minipage}[t]{\boxwidth}\footnotesize
%      \parskip=0pt\relax}%
%       {\end{minipage}\end{lrbox}%
%       \setlength{\xmcomlen}{\textwidth}%
%       \addtolength{\xmcomlen}{-\wd\alignbox}%
%       \settodepth{\alignwidth}{\usebox{\alignbox}}%
%       \global\setlength{\multicommentdepth}{\alignwidth}%
%       \setlength{\boxwidth}{\alignwidth}%
%       \global\addtolength{\alignwidth}{-\maxdepth}%
%       \addtolength{\boxwidth}{.1em}%
%       \raisebox{0pt}[0pt][0pt]{%
%        \ifthenelse{\boolean{shading}}%
%          {\hspace*{-\xmcomlen}\shadebox{\rule[-\boxwidth]{0pt}{0pt}%
%                                 \hspace*{\xmcomlen}\usebox{\alignbox}}}%
%          {\usebox{\alignbox}}}%
%       \vspace*{\alignwidth}\pagebreak[0]\vspace{-\alignwidth}\par}
% % a multi-line comment, whose first argument is its width in points.
%  
   
%%%%%%%%%%%%%%%%%%%%%%%%%%%%%%%%%%%%%%%%%%%%%%%%%%%%%%%%%%%%%%%%%%%%%%%%%%%%%%
% THE lcom ENVIRONMENT                                                       %
% ^^^^^^^^^^^^^^^^^^^^                                                       %
% A multi-line comment with no text to its left is typeset in an lcom        % 
% environment, whose argument is a number representing the indentation       % 
% of the left margin, in points.  All the text of the comment should be      % 
% inside cpar environments.                                                  % 
%%%%%%%%%%%%%%%%%%%%%%%%%%%%%%%%%%%%%%%%%%%%%%%%%%%%%%%%%%%%%%%%%%%%%%%%%%%%%%
\newenvironment{lcom}[1]{%
  \setlength{\lcomindent}{#1pt} % Added for PlusCal handling.
  \par\vspace{.2em}%
  \sloppypar
  \setcounter{pardepth}{0}%
  \footnotesize
  \begin{list}{}{%
    \setlength{\leftmargin}{#1pt}
    \setlength{\labelwidth}{0pt}%
    \setlength{\labelsep}{0pt}%
    \setlength{\itemindent}{0pt}%
    \setlength{\topsep}{0pt}%
    \setlength{\parsep}{0pt}%
    \setlength{\partopsep}{0pt}%
    \setlength{\parskip}{0pt}}
    \item[]}%
  {\end{list}\vspace{.3em}\setlength{\lcomindent}{0pt}%
 }


%%%%%%%%%%%%%%%%%%%%%%%%%%%%%%%%%%%%%%%%%%%%%%%%%%%%%%%%%%%%%%%%%%%%%%%%%%%%%
% THE mcom ENVIRONMENT AND \mutivspace COMMAND                              %
% ^^^^^^^^^^^^^^^^^^^^^^^^^^^^^^^^^^^^^^^^^^^^                              %
%                                                                           %
% A part of the spec containing a right-comment of the form                 %
%                                                                           %
%      xxxx (*************)                                                 %
%      yyyy (* ccccccccc *)                                                 %
%      ...  (* ccccccccc *)                                                 %
%           (* ccccccccc *)                                                 %
%           (* ccccccccc *)                                                 %
%           (*************)                                                 %
%                                                                           %
% is typeset by                                                             %
%                                                                           %
%     XXXX \begin{mcom}{d}                                                  %
%            CCCC ... CCC                                                   %
%          \end{mcom}                                                       %
%     YYYY ...                                                              %
%     \multivspace{n}                                                       %
%                                                                           %
% where the number d is the width in points of the comment, n is the        %
% number of xxxx, yyyy, ...  lines to the left of the comment.              %
% All the text of the comment should be typeset in cpar environments.       %
%                                                                           %
% This puts the comment into a single box (so no page breaks can occur      %
% within it).  The entire box is shaded iff the shading flag is true.       %
%%%%%%%%%%%%%%%%%%%%%%%%%%%%%%%%%%%%%%%%%%%%%%%%%%%%%%%%%%%%%%%%%%%%%%%%%%%%%
\newlength{\xmcomlen}%
\newenvironment{mcom}[1]{%
  \setcounter{pardepth}{0}%
  \hspace{.65em}%
  \begin{lrbox}{\alignbox}\sloppypar%
      \setboolean{shading}{false}%
      \setlength{\boxwidth}{#1pt}%
      \addtolength{\boxwidth}{-.65em}%
      \begin{minipage}[t]{\boxwidth}\footnotesize
      \parskip=0pt\relax}%
       {\end{minipage}\end{lrbox}%
       \setlength{\xmcomlen}{\textwidth}%       % For PlusCal shading
       \addtolength{\xmcomlen}{-\wd\alignbox}%  % For PlusCal shading
       \settodepth{\alignwidth}{\usebox{\alignbox}}%
       \global\setlength{\multicommentdepth}{\alignwidth}%
       \setlength{\boxwidth}{\alignwidth}%      % For PlusCal shading
       \global\addtolength{\alignwidth}{-\maxdepth}%
       \addtolength{\boxwidth}{.1em}%           % For PlusCal shading
      \raisebox{0pt}[0pt][0pt]{%
        \ifthenelse{\boolean{shading}}%
          {\ifpcalshading
             \hspace*{-\xmcomlen}%
             \shadebox{\rule[-\boxwidth]{0pt}{0pt}\hspace*{\xmcomlen}%
                          \usebox{\alignbox}}%
           \else
             \shadebox{\usebox{\alignbox}}
           \fi
          }%
          {\usebox{\alignbox}}}%
       \vspace*{\alignwidth}\pagebreak[0]\vspace{-\alignwidth}\par}
 % a multi-line comment, whose first argument is its width in points.


% \multispace{n} produces the vertical space indicated by "|"s in 
% this situation
%   
%     xxxx (*************)
%     xxxx (* ccccccccc *)
%      |   (* ccccccccc *)
%      |   (* ccccccccc *)
%      |   (* ccccccccc *)
%      |   (*************)
%
% where n is the number of "xxxx" lines.
\newcommand{\multivspace}[1]{\addtolength{\multicommentdepth}{-#1\baselineskip}%
 \addtolength{\multicommentdepth}{1.2em}%
 \ifthenelse{\lengthtest{\multicommentdepth > 0pt}}%
    {\par\vspace{\multicommentdepth}\par}{}}

%\newenvironment{hpar}[2]{%
%  \begin{list}{}{\setlength{\leftmargin}{#1pt}%
%                 \addtolength{\leftmargin}{#2pt}%
%                 \setlength{\itemindent}{-#2pt}%
%                 \setlength{\topsep}{0pt}%
%                 \setlength{\parsep}{0pt}%
%                 \setlength{\partopsep}{0pt}%
%                 \setlength{\parskip}{0pt}%
%                 \addtolength{\labelsep}{0pt}}%
%  \item[]\footnotesize}{\end{list}}
%    %%%%%%%%%%%%%%%%%%%%%%%%%%%%%%%%%%%%%%%%%%%%%%%%%%%%%%%%%%%%%%%%%%%%%%%%
%    % Typesets a sequence of paragraphs like this:                         %
%    %                                                                      %
%    %      left |<-- d1 --> XXXXXXXXXXXXXXXXXXXXXXXX                       %
%    %    margin |           <- d2 -> XXXXXXXXXXXXXXX                       %
%    %           |                    XXXXXXXXXXXXXXX                       %
%    %           |                                                          %
%    %           |                    XXXXXXXXXXXXXXX                       %
%    %           |                    XXXXXXXXXXXXXXX                       %
%    %                                                                      %
%    % where d1 = #1pt and d2 = #2pt, but with no vspace between            %
%    % paragraphs.                                                          %
%    %%%%%%%%%%%%%%%%%%%%%%%%%%%%%%%%%%%%%%%%%%%%%%%%%%%%%%%%%%%%%%%%%%%%%%%%

%%%%%%%%%%%%%%%%%%%%%%%%%%%%%%%%%%%%%%%%%%%%%%%%%%%%%%%%%%%%%%%%%%%%%%
% Commands for repeated characters that produce dashes.              %
%%%%%%%%%%%%%%%%%%%%%%%%%%%%%%%%%%%%%%%%%%%%%%%%%%%%%%%%%%%%%%%%%%%%%%
% \raisedDash{wd}{ht}{thk} makes a horizontal line wd characters wide, 
% raised a distance ht ex's above the baseline, with a thickness of 
% thk em's.
\newcommand{\raisedDash}[3]{\raisebox{#2ex}{\setlength{\alignwidth}{.5em}%
  \rule{#1\alignwidth}{#3em}}}

% The following commands take a single argument n and produce the
% output for n repeated characters, as follows
%   \cdash:    -
%   \tdash:    ~
%   \ceqdash:  =
%   \usdash:   _
\newcommand{\cdash}[1]{\raisedDash{#1}{.5}{.04}}
\newcommand{\usdash}[1]{\raisedDash{#1}{0}{.04}}
\newcommand{\ceqdash}[1]{\raisedDash{#1}{.5}{.08}}
\newcommand{\tdash}[1]{\raisedDash{#1}{1}{.08}}

\newlength{\spacewidth}
\setlength{\spacewidth}{.2em}
\newcommand{\e}[1]{\hspace{#1\spacewidth}}
%% \e{i} produces space corresponding to i input spaces.


%% Alignment-file Commands

\newlength{\alignboxwidth}
\newlength{\alignwidth}
\newsavebox{\alignbox}

% \al{i}{j}{txt} is used in the alignment file to put "%{i}{j}{wd}"
% in the log file, where wd is the width of the line up to that point,
% and txt is the following text.
\newcommand{\al}[3]{%
  \typeout{\%{#1}{#2}{\the\alignwidth}}%
  \cl{#3}}

%% \cl{txt} continues a specification line in the alignment file
%% with text txt.
\newcommand{\cl}[1]{%
  \savebox{\alignbox}{\mbox{$\mbox{}#1\mbox{}$}}%
  \settowidth{\alignboxwidth}{\usebox{\alignbox}}%
  \addtolength{\alignwidth}{\alignboxwidth}%
  \usebox{\alignbox}}

% \fl{txt} in the alignment file begins a specification line that
% starts with the text txt.
\newcommand{\fl}[1]{%
  \par
  \savebox{\alignbox}{\mbox{$\mbox{}#1\mbox{}$}}%
  \settowidth{\alignwidth}{\usebox{\alignbox}}%
  \usebox{\alignbox}}



  
%%%%%%%%%%%%%%%%%%%%%%%%%%%%%%%%%%%%%%%%%%%%%%%%%%%%%%%%%%%%%%%%%%%%%%%%%%%%%
% Ordinarily, TeX typesets letters in math mode in a special math italic    %
% font.  This makes it typeset "it" to look like the product of the         %
% variables i and t, rather than like the word "it".  The following         %
% commands tell TeX to use an ordinary italic font instead.                 %
%%%%%%%%%%%%%%%%%%%%%%%%%%%%%%%%%%%%%%%%%%%%%%%%%%%%%%%%%%%%%%%%%%%%%%%%%%%%%
\ifx\documentclass\undefined
\else
  \DeclareSymbolFont{tlaitalics}{\encodingdefault}{cmr}{m}{it}
  \let\itfam\symtlaitalics
\fi

\makeatletter
\newcommand{\tlx@c}{\c@tlx@ctr\advance\c@tlx@ctr\@ne}
\newcounter{tlx@ctr}
\c@tlx@ctr=\itfam \multiply\c@tlx@ctr"100\relax \advance\c@tlx@ctr "7061\relax
\mathcode`a=\tlx@c \mathcode`b=\tlx@c \mathcode`c=\tlx@c \mathcode`d=\tlx@c
\mathcode`e=\tlx@c \mathcode`f=\tlx@c \mathcode`g=\tlx@c \mathcode`h=\tlx@c
\mathcode`i=\tlx@c \mathcode`j=\tlx@c \mathcode`k=\tlx@c \mathcode`l=\tlx@c
\mathcode`m=\tlx@c \mathcode`n=\tlx@c \mathcode`o=\tlx@c \mathcode`p=\tlx@c
\mathcode`q=\tlx@c \mathcode`r=\tlx@c \mathcode`s=\tlx@c \mathcode`t=\tlx@c
\mathcode`u=\tlx@c \mathcode`v=\tlx@c \mathcode`w=\tlx@c \mathcode`x=\tlx@c
\mathcode`y=\tlx@c \mathcode`z=\tlx@c
\c@tlx@ctr=\itfam \multiply\c@tlx@ctr"100\relax \advance\c@tlx@ctr "7041\relax
\mathcode`A=\tlx@c \mathcode`B=\tlx@c \mathcode`C=\tlx@c \mathcode`D=\tlx@c
\mathcode`E=\tlx@c \mathcode`F=\tlx@c \mathcode`G=\tlx@c \mathcode`H=\tlx@c
\mathcode`I=\tlx@c \mathcode`J=\tlx@c \mathcode`K=\tlx@c \mathcode`L=\tlx@c
\mathcode`M=\tlx@c \mathcode`N=\tlx@c \mathcode`O=\tlx@c \mathcode`P=\tlx@c
\mathcode`Q=\tlx@c \mathcode`R=\tlx@c \mathcode`S=\tlx@c \mathcode`T=\tlx@c
\mathcode`U=\tlx@c \mathcode`V=\tlx@c \mathcode`W=\tlx@c \mathcode`X=\tlx@c
\mathcode`Y=\tlx@c \mathcode`Z=\tlx@c
\makeatother

%%%%%%%%%%%%%%%%%%%%%%%%%%%%%%%%%%%%%%%%%%%%%%%%%%%%%%%%%%
%                THE describe ENVIRONMENT                %
%%%%%%%%%%%%%%%%%%%%%%%%%%%%%%%%%%%%%%%%%%%%%%%%%%%%%%%%%%
%
%
% It is like the description environment except it takes an argument
% ARG that should be the text of the widest label.  It adjusts the
% indentation so each item with label LABEL produces
%%      LABEL             blah blah blah
%%      <- width of ARG ->blah blah blah
%%                        blah blah blah
\newenvironment{describe}[1]%
   {\begin{list}{}{\settowidth{\labelwidth}{#1}%
            \setlength{\labelsep}{.5em}%
            \setlength{\leftmargin}{\labelwidth}% 
            \addtolength{\leftmargin}{\labelsep}%
            \addtolength{\leftmargin}{\parindent}%
            \def\makelabel##1{\rm ##1\hfill}}%
            \setlength{\topsep}{0pt}}%% 
                % Sets \topsep to 0 to reduce vertical space above
                % and below embedded displayed equations
   {\end{list}}

%   For tlatex.TeX
\usepackage{verbatim}
\makeatletter
\def\tla{\let\%\relax%
         \@bsphack
         \typeout{\%{\the\linewidth}}%
             \let\do\@makeother\dospecials\catcode`\^^M\active
             \let\verbatim@startline\relax
             \let\verbatim@addtoline\@gobble
             \let\verbatim@processline\relax
             \let\verbatim@finish\relax
             \verbatim@}
\let\endtla=\@esphack

\let\pcal=\tla
\let\endpcal=\endtla
\let\ppcal=\tla
\let\endppcal=\endtla

% The tlatex environment is used by TLATeX.TeX to typeset TLA+.
% TLATeX.TLA starts its files by writing a \tlatex command.  This
% command/environment sets \parindent to 0 and defines \% to its
% standard definition because the writing of the log files is messed up
% if \% is defined to be something else.  It also executes
% \@computerule to determine the dimensions for the TLA horizonatl
% bars.
\newenvironment{tlatex}{\@computerule%
                        \setlength{\parindent}{0pt}%
                       \makeatletter\chardef\%=`\%}{}


% The notla environment produces no output.  You can turn a 
% tla environment to a notla environment to prevent tlatex.TeX from
% re-formatting the environment.

\def\notla{\let\%\relax%
         \@bsphack
             \let\do\@makeother\dospecials\catcode`\^^M\active
             \let\verbatim@startline\relax
             \let\verbatim@addtoline\@gobble
             \let\verbatim@processline\relax
             \let\verbatim@finish\relax
             \verbatim@}
\let\endnotla=\@esphack

\let\nopcal=\notla
\let\endnopcal=\endnotla
\let\noppcal=\notla
\let\endnoppcal=\endnotla

%%%%%%%%%%%%%%%%%%%%%%%% end of tlatex.sty file %%%%%%%%%%%%%%%%%%%%%%% 
% last modified on Fri  3 August 2012 at 14:23:49 PST by lamport

\begin{document}
\tlatex
 \@x{\makebox[0pt][r]{\scriptsize 1\hspace{1em}}}\moduleLeftDash\@xx{
 {\MODULE} Overwatch}\moduleRightDash\@xx{}%
\begin{lcom}{0}%
\begin{cpar}{0}{T}{F}{7.5}{0}{}%
Section: Imports
\end{cpar}%
\vshade{5.0}%
\begin{cpar}{1}{F}{F}{0}{0}{}%
Shared operators used by this module are imported below.
\end{cpar}%
\end{lcom}%
 \@x{\makebox[0pt][r]{\scriptsize 6\hspace{1em}} {\EXTENDS} Naturals ,\,
 FiniteSets ,\, Sequences ,\, FailurePolicy}%
\@pvspace{8.0pt}%
\begin{lcom}{0}%
\begin{cpar}{0}{T}{F}{7.5}{0}{}%
Section: State and result domains
\end{cpar}%
\vshade{5.0}%
\begin{cpar}{1}{F}{F}{0}{0}{}%
 Workspace versions, completed outcomes, and message records form the result
 layer.
 Finite bounds are supplied by the \ensuremath{TLC} configuration, not by the
 runtime.
\end{cpar}%
\end{lcom}%
 \@x{\makebox[0pt][r]{\scriptsize 13\hspace{1em}} {\CONSTANTS} Names ,\,
 MaxEdits ,\, Timeout}%
 \@x{\makebox[0pt][r]{\scriptsize 14\hspace{1em}} {\ASSUME}\@w{timeout}
 \.{\notin} Names}%
\@pvspace{8.0pt}%
 \@x{\makebox[0pt][r]{\scriptsize 16\hspace{1em}} Versions \.{\defeq} 0
 \.{\dotdot} MaxEdits}%
 \@x{\makebox[0pt][r]{\scriptsize 17\hspace{1em}} Results \.{\defeq} {\SUBSET}
 Names}%
 \@x{\makebox[0pt][r]{\scriptsize 18\hspace{1em}} TimedOut
 \.{\defeq}\@w{timeout}}%
 \@x{\makebox[0pt][r]{\scriptsize 19\hspace{1em}} Outcomes\@s{1.66} \.{\defeq}
 {\SUBSET} ( Names \.{\cup} \{ TimedOut \} )}%
 \@x{\makebox[0pt][r]{\scriptsize 20\hspace{1em}} NoMessage \.{\defeq} [ held
 \.{\mapsto} {\FALSE} ,\, waking \.{\mapsto} {\FALSE} ,\, failures
 \.{\mapsto} \{ \} ]}%
 \@x{\makebox[0pt][r]{\scriptsize 21\hspace{1em}} NotTold \.{\defeq} [ kind
 \.{\mapsto}\@w{none} ,\, fresh\@s{6.99} \.{\mapsto} \{ \} ]}%
 \@x{\makebox[0pt][r]{\scriptsize 22\hspace{1em}} NoBaseline \.{\defeq} [
 known \.{\mapsto} {\FALSE} ,\, at \.{\mapsto} 0 ,\, failures \.{\mapsto} \{
 \} ]}%
\@pvspace{8.0pt}%
\begin{lcom}{0}%
\begin{cpar}{0}{T}{F}{7.5}{0}{}%
Section: State shape
\end{cpar}%
\vshade{5.0}%
\begin{cpar}{1}{F}{F}{0}{0}{}%
 The checked version, unread results, inherited baseline, and delivery state
 are
 separate. touched records whether the session has edited since the commit the
 baseline belongs to; own is the versions those edits made, which only the
 properties read. \ensuremath{TypeOK} constrains their domains without
 advancing an execution.
\end{cpar}%
\end{lcom}%
\@x{\makebox[0pt][r]{\scriptsize 31\hspace{1em}} {\VARIABLES}}%
 \@x{\makebox[0pt][r]{\scriptsize 32\hspace{1em}}\@s{16.4} truth ,\, head ,\,
 based ,\, busy ,\, from ,\, quiet ,\, installed ,\,}%
 \@x{\makebox[0pt][r]{\scriptsize 33\hspace{1em}}\@s{16.4} produced ,\,
 consumed ,\, pending ,\, report ,\, baseline ,\, wasFailing ,\,}%
 \@x{\makebox[0pt][r]{\scriptsize 34\hspace{1em}}\@s{16.4} outbox ,\, told ,\,
 flush ,\, turn ,\, toldThisTurn ,\, promptQueued ,\, knows ,\, elapsed ,\,
 edited ,\,}%
\@x{\makebox[0pt][r]{\scriptsize 35\hspace{1em}}\@s{16.4} touched ,\, own}%
\@pvspace{8.0pt}%
 \@x{\makebox[0pt][r]{\scriptsize 37\hspace{1em}} vars \.{\defeq} {\langle}
 truth ,\, head ,\, based ,\, busy ,\, from ,\, quiet ,\, installed ,\,}%
 \@x{\makebox[0pt][r]{\scriptsize 38\hspace{1em}}\@s{41.61} produced ,\,
 consumed ,\, pending ,\, report ,\, baseline ,\, wasFailing ,\,}%
 \@x{\makebox[0pt][r]{\scriptsize 39\hspace{1em}}\@s{41.61} outbox ,\, told
 ,\, flush ,\, turn ,\, toldThisTurn ,\, promptQueued ,\, knows ,\, elapsed
 ,\, edited ,\,}%
 \@x{\makebox[0pt][r]{\scriptsize 40\hspace{1em}}\@s{41.61} touched ,\, own
 {\rangle}}%
\@pvspace{8.0pt}%
\@x{\makebox[0pt][r]{\scriptsize 42\hspace{1em}} TypeOK \.{\defeq}}%
 \@x{\makebox[0pt][r]{\scriptsize 43\hspace{1em}}\@s{16.4} \.{\land}\@s{9.74}
 truth \.{\in} [ Versions \.{\rightarrow} Results ]}%
 \@x{\makebox[0pt][r]{\scriptsize 44\hspace{1em}}\@s{16.4} \.{\land}\@s{9.74}
 head \.{\in} Versions \.{\land} based \.{\in} 0 \.{\dotdot} head}%
 \@x{\makebox[0pt][r]{\scriptsize 45\hspace{1em}}\@s{16.4} \.{\land}\@s{9.74}
 busy\@s{0.65} \.{\in} {\BOOLEAN} \.{\land} from \.{\in} Versions \.{\land}
 quiet \.{\in} {\BOOLEAN} \.{\land} installed \.{\in} {\BOOLEAN}}%
 \@x{\makebox[0pt][r]{\scriptsize 46\hspace{1em}}\@s{16.4} \.{\land}\@s{9.74}
 produced \.{\in} 0 \.{\dotdot} ( MaxEdits \.{+} 1 ) \.{\land} consumed
 \.{\in} 0 \.{\dotdot} produced}%
 \@x{\makebox[0pt][r]{\scriptsize 47\hspace{1em}}\@s{16.4} \.{\land}\@s{9.74}
 pending \.{\in} Seq ( [ id \.{:} 1 \.{\dotdot} ( MaxEdits \.{+} 1 ) ,\, at
 \.{:} Versions ,\, failures \.{:} Outcomes ] )}%
 \@x{\makebox[0pt][r]{\scriptsize 48\hspace{1em}}\@s{16.4} \.{\land}\@s{9.74}
 Len ( pending ) \.{\leq} MaxEdits \.{+} 1}%
 \@x{\makebox[0pt][r]{\scriptsize 49\hspace{1em}}\@s{16.4} \.{\land}\@s{9.74}
 report \.{\in} [ known \.{:} {\BOOLEAN} ,\, id \.{:} 0 \.{\dotdot} produced
 ,\, at \.{:} Versions ,\, failures \.{:} Outcomes ]}%
 \@x{\makebox[0pt][r]{\scriptsize 50\hspace{1em}}\@s{16.4} \.{\land}\@s{9.74}
 baseline \.{\in} [ known \.{:} {\BOOLEAN} ,\, at \.{:} Versions ,\, failures
 \.{:} Outcomes ]}%
 \@x{\makebox[0pt][r]{\scriptsize 51\hspace{1em}}\@s{16.4} \.{\land}\@s{9.74}
 wasFailing \.{\in} {\BOOLEAN}}%
 \@x{\makebox[0pt][r]{\scriptsize 52\hspace{1em}}\@s{16.4} \.{\land}\@s{9.74}
 outbox \.{\in} [ held \.{:} {\BOOLEAN} ,\, waking \.{:} {\BOOLEAN} ,\,
 failures \.{:} Outcomes ]}%
 \@x{\makebox[0pt][r]{\scriptsize 53\hspace{1em}}\@s{16.4} \.{\land}\@s{9.74}
 told \.{\in} [ kind \.{:} \{\@w{none} ,\,\@w{failed} ,\,\@w{recovered} \}
 ,\, fresh \.{:} Outcomes ]}%
 \@x{\makebox[0pt][r]{\scriptsize 54\hspace{1em}}\@s{16.4} \.{\land}\@s{9.74}
 flush \.{\in} {\BOOLEAN}}%
 \@x{\makebox[0pt][r]{\scriptsize 55\hspace{1em}}\@s{16.4} \.{\land}\@s{9.74}
 turn \.{\in} \{\@w{idle} ,\,\@w{running} \}}%
 \@x{\makebox[0pt][r]{\scriptsize 56\hspace{1em}}\@s{16.4} \.{\land}\@s{9.74}
 toldThisTurn \.{\in} {\BOOLEAN} \.{\land} promptQueued \.{\in} {\BOOLEAN}}%
 \@x{\makebox[0pt][r]{\scriptsize 57\hspace{1em}}\@s{16.4} \.{\land}\@s{9.74}
 knows \.{\in} Outcomes \.{\land} elapsed \.{\in} 0 \.{\dotdot} Timeout}%
 \@x{\makebox[0pt][r]{\scriptsize 58\hspace{1em}}\@s{16.4} \.{\land}\@s{9.74}
 edited \.{\in} {\BOOLEAN}}%
 \@x{\makebox[0pt][r]{\scriptsize 59\hspace{1em}}\@s{16.4} \.{\land}\@s{9.74}
 touched \.{\in} {\BOOLEAN} \.{\land} own \.{\subseteq} Versions}%
\@pvspace{8.0pt}%
\begin{lcom}{0}%
\begin{cpar}{0}{T}{F}{7.5}{0}{}%
Section: Failure gating
\end{cpar}%
\vshade{5.0}%
\begin{cpar}{1}{F}{F}{0}{0}{}%
A failure report becomes a gate only when it is a regression against the
 session baseline, or when the check times out.
\end{cpar}%
\end{lcom}%
 \@x{\makebox[0pt][r]{\scriptsize 66\hspace{1em}} Failing \.{\defeq} report .
 known \.{\land} report . failures \.{\neq} \{ \}}%
 \@x{\makebox[0pt][r]{\scriptsize 67\hspace{1em}} Regressed ( f ) \.{\defeq}
 TimedOut \.{\in} f \.{\lor} ( baseline . known \.{\land} NewFailure ( f ,\,
 baseline . failures ) )}%
 \@x{\makebox[0pt][r]{\scriptsize 68\hspace{1em}} Guard \.{\defeq} Failing
 \.{\land} Regressed ( report . failures )}%
\@pvspace{8.0pt}%
\begin{lcom}{0}%
\begin{cpar}{0}{T}{F}{7.5}{0}{}%
Section: Initial state
\end{cpar}%
\vshade{5.0}%
\begin{cpar}{1}{F}{F}{0}{0}{}%
The first check starts at the inherited workspace version. No result or
 baseline is known and no notification has been sent.
\end{cpar}%
\end{lcom}%
\@x{\makebox[0pt][r]{\scriptsize 75\hspace{1em}} Init \.{\defeq}}%
 \@x{\makebox[0pt][r]{\scriptsize 76\hspace{1em}}\@s{16.4} \.{\land} truth
 \.{\in} [ Versions \.{\rightarrow} Results ]}%
 \@x{\makebox[0pt][r]{\scriptsize 77\hspace{1em}}\@s{16.4} \.{\land} head
 \.{=} 0 \.{\land} based \.{=} 0}%
 \@x{\makebox[0pt][r]{\scriptsize 78\hspace{1em}}\@s{16.4} \.{\land}
 busy\@s{0.65} \.{=} {\TRUE} \.{\land} from \.{=} 0 \.{\land} quiet \.{=}
 {\FALSE} \.{\land} installed \.{=} {\FALSE}}%
 \@x{\makebox[0pt][r]{\scriptsize 79\hspace{1em}}\@s{16.4} \.{\land} produced
 \.{=} 0 \.{\land} consumed \.{=} 0 \.{\land} pending \.{=} {\langle}
 {\rangle}}%
 \@x{\makebox[0pt][r]{\scriptsize 80\hspace{1em}}\@s{16.4} \.{\land} report
 \.{=} [ known \.{\mapsto} {\FALSE} ,\, id \.{\mapsto} 0 ,\, at \.{\mapsto} 0
 ,\, failures \.{\mapsto} \{ \} ]}%
 \@x{\makebox[0pt][r]{\scriptsize 81\hspace{1em}}\@s{16.4} \.{\land} baseline
 \.{=} NoBaseline}%
 \@x{\makebox[0pt][r]{\scriptsize 82\hspace{1em}}\@s{16.4} \.{\land}
 wasFailing \.{=} {\FALSE}}%
 \@x{\makebox[0pt][r]{\scriptsize 83\hspace{1em}}\@s{16.4} \.{\land} outbox
 \.{=} NoMessage \.{\land} told \.{=} NotTold \.{\land} flush \.{=} {\FALSE}}%
 \@x{\makebox[0pt][r]{\scriptsize 84\hspace{1em}}\@s{16.4} \.{\land} turn
 \.{=}\@w{idle} \.{\land} toldThisTurn \.{=} {\FALSE} \.{\land} promptQueued
 \.{=} {\FALSE}}%
 \@x{\makebox[0pt][r]{\scriptsize 85\hspace{1em}}\@s{16.4} \.{\land} knows
 \.{=} \{ \} \.{\land} elapsed \.{=} 0 \.{\land} edited \.{=} {\FALSE}}%
 \@x{\makebox[0pt][r]{\scriptsize 86\hspace{1em}}\@s{16.4} \.{\land} touched
 \.{=} {\FALSE} \.{\land} own \.{=} \{ \}}%
\@pvspace{8.0pt}%
\begin{lcom}{0}%
\begin{cpar}{0}{T}{F}{7.5}{0}{}%
Section: Result delivery and baseline isolation
\end{cpar}%
\vshade{5.0}%
\begin{cpar}{1}{F}{F}{0}{0}{}%
Queue deduplicates notification signatures. Deliver consumes a completed
 snapshot, establishes the inherited baseline, and queues a regression or
 recovery.
 The first result is the baseline unless the session has edited since the
 commit, as another process\mbox{'}s changes are inherited like the commit
 itself; after
 a session edit the baseline is empty, so no failure the edit may hold is
 absorbed.
 A recovery can cancel a failure message that has not yet been delivered.
\end{cpar}%
\end{lcom}%
 \@x{\makebox[0pt][r]{\scriptsize 97\hspace{1em}} Queue ( signature ,\,
 message ) \.{\defeq}}%
 \@x{\makebox[0pt][r]{\scriptsize 98\hspace{1em}}\@s{16.4} {\IF} told \.{=}
 signature \.{\land} {\lnot} outbox . held}%
 \@x{\makebox[0pt][r]{\scriptsize 99\hspace{1em}}\@s{16.4} \.{\THEN}
 {\UNCHANGED} {\langle} outbox ,\, told ,\, flush {\rangle}}%
 \@x{\makebox[0pt][r]{\scriptsize 100\hspace{1em}}\@s{16.4} \.{\ELSE}
 \.{\land} told \.{'} \.{=} signature \.{\land} outbox \.{'} \.{=} message
 \.{\land} flush \.{'} \.{=} {\TRUE}}%
\@pvspace{8.0pt}%
 \@x{\makebox[0pt][r]{\scriptsize 102\hspace{1em}} Deliver ( v ,\, id ,\, f )
 \.{\defeq}}%
 \@x{\makebox[0pt][r]{\scriptsize 103\hspace{1em}}\@s{16.4} \.{\LET} b
 \.{\defeq} {\IF} baseline . known \.{\THEN} baseline . failures \.{\ELSE}
 {\IF} {\lnot} touched \.{\THEN} f \.{\ELSE} \{ \}}%
 \@x{\makebox[0pt][r]{\scriptsize 104\hspace{1em}}\@s{36.79} isRegression
 \.{\defeq} TimedOut \.{\in} f \.{\lor} {\lnot} ( f \.{\subseteq} b )}%
 \@x{\makebox[0pt][r]{\scriptsize 105\hspace{1em}}\@s{16.4} \.{\IN} \.{\land}
 report \.{'} \.{=} [ known \.{\mapsto} {\TRUE} ,\, id \.{\mapsto} id ,\, at
 \.{\mapsto} v ,\, failures \.{\mapsto} f ]}%
 \@x{\makebox[0pt][r]{\scriptsize 106\hspace{1em}}\@s{36.79} \.{\land}
 baseline \.{'} \.{=} {\IF} TimedOut \.{\in} f \.{\lor} baseline . known
 \.{\THEN} baseline}%
 \@x{\makebox[0pt][r]{\scriptsize 107\hspace{1em}}\@s{98.53} \.{\ELSE} [ known
 \.{\mapsto} {\TRUE} ,\, at \.{\mapsto} v ,\, failures \.{\mapsto} b ]}%
 \@x{\makebox[0pt][r]{\scriptsize 108\hspace{1em}}\@s{36.79} \.{\land}
 wasFailing \.{'} \.{=} isRegression}%
 \@x{\makebox[0pt][r]{\scriptsize 109\hspace{1em}}\@s{36.79} \.{\land} {\IF}
 isRegression}%
 \@x{\makebox[0pt][r]{\scriptsize 110\hspace{1em}}\@s{56.11} \.{\THEN} Queue (
 [ kind \.{\mapsto}\@w{failed} ,\, fresh \.{\mapsto} f \.{\,\backslash\,} b ]
 ,\,}%
 \@x{\makebox[0pt][r]{\scriptsize 111\hspace{1em}}\@s{119.66} [ held\@s{1.02}
 \.{\mapsto} {\TRUE} ,\, waking \.{\mapsto} edited \.{\lor} TimedOut \.{\in}
 f ,\, failures \.{\mapsto} f ] )}%
 \@x{\makebox[0pt][r]{\scriptsize 112\hspace{1em}}\@s{47.91} \.{\ELSE} {\IF}
 {\lnot} wasFailing}%
 \@x{\makebox[0pt][r]{\scriptsize 113\hspace{1em}}\@s{56.11} \.{\THEN}
 {\UNCHANGED} {\langle} outbox ,\, told ,\, flush {\rangle}}%
 \@x{\makebox[0pt][r]{\scriptsize 114\hspace{1em}}\@s{47.91} \.{\ELSE} {\IF}
 outbox . held}%
 \@x{\makebox[0pt][r]{\scriptsize 115\hspace{1em}}\@s{56.11} \.{\THEN}
 \.{\land} outbox \.{'} \.{=} NoMessage \.{\land} told \.{'} \.{=} NotTold
 \.{\land} {\UNCHANGED} flush}%
 \@x{\makebox[0pt][r]{\scriptsize 116\hspace{1em}}\@s{47.91} \.{\ELSE} Queue (
 [ kind \.{\mapsto}\@w{recovered} ,\, fresh \.{\mapsto} \{ \} ] ,\, [ held
 \.{\mapsto} {\TRUE} ,\, waking \.{\mapsto} {\FALSE} ,\, failures \.{\mapsto}
 f ] )}%
\@pvspace{8.0pt}%
\begin{lcom}{0}%
\begin{cpar}{0}{T}{F}{7.5}{0}{}%
Section: Workspace changes and scheduling
\end{cpar}%
\vshade{5.0}%
\begin{cpar}{1}{F}{F}{0}{0}{}%
An agent edit advances head without moving based, and marks the next read as
 following the agent\mbox{'}s own change and the session as having edited.
 Another
 process can change the workspace at any time without committing;
 \ensuremath{OutsideEdit
 } advances head the same way but marks nothing. An external head change moves
 both,
 invalidates the baseline and unread results, and starts the session\mbox{'}s
 edits
 afresh. \ensuremath{QuietElapses} starts the next check; \ensuremath{Tick}
 advances its timeout clock.
\end{cpar}%
\end{lcom}%
\@x{\makebox[0pt][r]{\scriptsize 127\hspace{1em}} Edit \.{\defeq}}%
 \@x{\makebox[0pt][r]{\scriptsize 128\hspace{1em}}\@s{16.4} \.{\land} turn
 \.{=}\@w{running}\@s{0.17} \.{\land} head \.{<} MaxEdits}%
 \@x{\makebox[0pt][r]{\scriptsize 129\hspace{1em}}\@s{16.4} \.{\land} head
 \.{'} \.{=} head \.{+} 1 \.{\land} quiet \.{'} \.{=} {\TRUE} \.{\land}
 edited \.{'} \.{=} {\TRUE}}%
 \@x{\makebox[0pt][r]{\scriptsize 130\hspace{1em}}\@s{16.4} \.{\land} touched
 \.{'} \.{=} {\TRUE} \.{\land} own \.{'} \.{=} own \.{\cup} \{ head \.{+} 1
 \}}%
 \@x{\makebox[0pt][r]{\scriptsize 131\hspace{1em}}\@s{16.4} \.{\land}
 {\UNCHANGED} {\langle} truth ,\, based ,\, busy ,\, from ,\, installed ,\,
 produced ,\, consumed ,\, pending ,\,}%
 \@x{\makebox[0pt][r]{\scriptsize 132\hspace{1em}}\@s{89.98} report ,\,
 baseline ,\, wasFailing ,\, outbox ,\, told ,\, flush ,\, turn ,\,}%
 \@x{\makebox[0pt][r]{\scriptsize 133\hspace{1em}}\@s{89.98} toldThisTurn ,\,
 promptQueued ,\, knows ,\, elapsed {\rangle}}%
\@pvspace{8.0pt}%
\@x{\makebox[0pt][r]{\scriptsize 135\hspace{1em}} OutsideEdit \.{\defeq}}%
 \@x{\makebox[0pt][r]{\scriptsize 136\hspace{1em}}\@s{16.4} \.{\land} head
 \.{<} MaxEdits}%
 \@x{\makebox[0pt][r]{\scriptsize 137\hspace{1em}}\@s{16.4} \.{\land} head
 \.{'} \.{=} head \.{+} 1 \.{\land} quiet \.{'} \.{=} {\TRUE}}%
 \@x{\makebox[0pt][r]{\scriptsize 138\hspace{1em}}\@s{16.4} \.{\land}
 {\UNCHANGED} {\langle} truth ,\, based ,\, busy ,\, from ,\, installed ,\,
 produced ,\, consumed ,\, pending ,\,}%
 \@x{\makebox[0pt][r]{\scriptsize 139\hspace{1em}}\@s{89.98} report ,\,
 baseline ,\, wasFailing ,\, outbox ,\, told ,\, flush ,\, turn ,\,}%
 \@x{\makebox[0pt][r]{\scriptsize 140\hspace{1em}}\@s{89.98} toldThisTurn ,\,
 promptQueued ,\, knows ,\, elapsed ,\, edited ,\, touched ,\, own {\rangle}}%
\@pvspace{8.0pt}%
\@x{\makebox[0pt][r]{\scriptsize 142\hspace{1em}} HeadMoves \.{\defeq}}%
 \@x{\makebox[0pt][r]{\scriptsize 143\hspace{1em}}\@s{16.4} \.{\land} head
 \.{<} MaxEdits}%
 \@x{\makebox[0pt][r]{\scriptsize 144\hspace{1em}}\@s{16.4} \.{\land} head
 \.{'} \.{=} head \.{+} 1 \.{\land} based \.{'} \.{=} head \.{+} 1}%
 \@x{\makebox[0pt][r]{\scriptsize 145\hspace{1em}}\@s{16.4} \.{\land} baseline
 \.{'} \.{=} NoBaseline \.{\land} touched \.{'} \.{=} {\FALSE} \.{\land} own
 \.{'} \.{=} \{ \}}%
 \@x{\makebox[0pt][r]{\scriptsize 146\hspace{1em}}\@s{16.4} \.{\land} quiet
 \.{'} \.{=} {\TRUE} \.{\land} pending \.{'} \.{=} {\langle} {\rangle}}%
 \@x{\makebox[0pt][r]{\scriptsize 147\hspace{1em}}\@s{16.4} \.{\land} from
 \.{'} \.{=} {\IF} busy \.{\THEN} head \.{+} 1 \.{\ELSE} from}%
 \@x{\makebox[0pt][r]{\scriptsize 148\hspace{1em}}\@s{16.4} \.{\land} elapsed
 \.{'} \.{=} {\IF} busy \.{\THEN} 0 \.{\ELSE} elapsed}%
 \@x{\makebox[0pt][r]{\scriptsize 149\hspace{1em}}\@s{16.4} \.{\land}
 {\UNCHANGED} {\langle} truth ,\, busy ,\, installed ,\, produced ,\,
 consumed ,\, report ,\, wasFailing ,\, outbox ,\,}%
 \@x{\makebox[0pt][r]{\scriptsize 150\hspace{1em}}\@s{89.98} told ,\, flush
 ,\, turn ,\, toldThisTurn ,\, promptQueued ,\, knows ,\, edited {\rangle}}%
\@pvspace{8.0pt}%
\@x{\makebox[0pt][r]{\scriptsize 152\hspace{1em}} QuietElapses \.{\defeq}}%
 \@x{\makebox[0pt][r]{\scriptsize 153\hspace{1em}}\@s{16.4} \.{\land} quiet
 \.{\land} installed}%
 \@x{\makebox[0pt][r]{\scriptsize 154\hspace{1em}}\@s{16.4} \.{\land} quiet
 \.{'} \.{=} {\FALSE} \.{\land} busy \.{'} \.{=} {\TRUE} \.{\land} from \.{'}
 \.{=} head \.{\land} elapsed \.{'} \.{=} 0}%
 \@x{\makebox[0pt][r]{\scriptsize 155\hspace{1em}}\@s{16.4} \.{\land}
 {\UNCHANGED} {\langle} truth ,\, head ,\, based ,\, installed ,\, produced
 ,\, consumed ,\, pending ,\,}%
 \@x{\makebox[0pt][r]{\scriptsize 156\hspace{1em}}\@s{89.98} report ,\,
 baseline ,\, wasFailing ,\, outbox ,\, told ,\, flush ,\, turn ,\,}%
 \@x{\makebox[0pt][r]{\scriptsize 157\hspace{1em}}\@s{89.98} toldThisTurn ,\,
 promptQueued ,\, knows ,\, edited ,\, touched ,\, own {\rangle}}%
\@pvspace{8.0pt}%
\@x{\makebox[0pt][r]{\scriptsize 159\hspace{1em}} Tick \.{\defeq}}%
 \@x{\makebox[0pt][r]{\scriptsize 160\hspace{1em}}\@s{16.4} \.{\land} busy
 \.{\land} elapsed \.{<} Timeout \.{\land} elapsed \.{'} \.{=} elapsed \.{+}
 1}%
 \@x{\makebox[0pt][r]{\scriptsize 161\hspace{1em}}\@s{16.4} \.{\land}
 {\UNCHANGED} {\langle} truth ,\, head ,\, based ,\, busy ,\, from ,\, quiet
 ,\, installed ,\, produced ,\,}%
 \@x{\makebox[0pt][r]{\scriptsize 162\hspace{1em}}\@s{89.98} consumed ,\,
 pending ,\, report ,\, baseline ,\, wasFailing ,\, outbox ,\,}%
 \@x{\makebox[0pt][r]{\scriptsize 163\hspace{1em}}\@s{89.98} told ,\, flush
 ,\, turn ,\, toldThisTurn ,\, promptQueued ,\, knows ,\, edited ,\,}%
 \@x{\makebox[0pt][r]{\scriptsize 164\hspace{1em}}\@s{89.98} touched ,\, own
 {\rangle}}%
\@pvspace{8.0pt}%
\begin{lcom}{0}%
\begin{cpar}{0}{T}{F}{7.5}{0}{}%
Completion appends a result; it does not deliver it. Partial output is a
 stuttering step at this abstraction: only a complete build enters pending.
\end{cpar}%
\begin{cpar}{1}{T}{F}{2.5}{0}{}%
Section: Buffered completion and consumption
\end{cpar}%
\vshade{5.0}%
\begin{cpar}{1}{F}{F}{0}{0}{}%
Run completion appends a snapshot to pending without publishing it.
 \ensuremath{ReadResults} consumes the newest completed snapshot and clears
 the buffer.
 A later incomplete build does not replace an earlier complete result.
\end{cpar}%
\end{lcom}%
 \@x{\makebox[0pt][r]{\scriptsize 174\hspace{1em}} Buffer ( v ,\, f )
 \.{\defeq}}%
 \@x{\makebox[0pt][r]{\scriptsize 175\hspace{1em}}\@s{16.4} \.{\land} produced
 \.{'} \.{=} produced \.{+} 1}%
 \@x{\makebox[0pt][r]{\scriptsize 176\hspace{1em}}\@s{16.4} \.{\land} pending
 \.{'} \.{=} Append ( pending ,\, [ id \.{\mapsto} produced \.{+} 1 ,\, at
 \.{\mapsto} v ,\, failures \.{\mapsto} f ] )}%
 \@x{\makebox[0pt][r]{\scriptsize 177\hspace{1em}}\@s{16.4} \.{\land} busy
 \.{'} \.{=} {\FALSE} \.{\land} installed \.{'} \.{=} {\TRUE}}%
 \@x{\makebox[0pt][r]{\scriptsize 178\hspace{1em}}\@s{16.4} \.{\land}
 {\UNCHANGED} {\langle} truth ,\, head ,\, based ,\, from ,\, quiet ,\,
 consumed ,\, report ,\, baseline ,\,}%
 \@x{\makebox[0pt][r]{\scriptsize 179\hspace{1em}}\@s{89.98} wasFailing ,\,
 outbox ,\, told ,\, flush ,\, turn ,\, toldThisTurn ,\, promptQueued ,\,
 knows ,\, elapsed ,\, edited ,\,}%
 \@x{\makebox[0pt][r]{\scriptsize 180\hspace{1em}}\@s{89.98} touched ,\, own
 {\rangle}}%
\@pvspace{8.0pt}%
 \@x{\makebox[0pt][r]{\scriptsize 182\hspace{1em}} RunEnds \.{\defeq} busy
 \.{\land} \E\, v \.{\in} from \.{\dotdot} head \.{:} Buffer ( v ,\, truth [
 v ] )}%
 \@x{\makebox[0pt][r]{\scriptsize 183\hspace{1em}} TimesOut \.{\defeq} busy
 \.{\land} elapsed \.{=} Timeout \.{\land} Buffer ( from ,\, \{ TimedOut \}
 )}%
\@pvspace{8.0pt}%
\@x{\makebox[0pt][r]{\scriptsize 185\hspace{1em}} ReadResults \.{\defeq}}%
 \@x{\makebox[0pt][r]{\scriptsize 186\hspace{1em}}\@s{16.4} \.{\land} pending
 \.{\neq} {\langle} {\rangle}}%
 \@x{\makebox[0pt][r]{\scriptsize 187\hspace{1em}}\@s{16.4} \.{\land} \.{\LET}
 latest \.{\defeq} pending [ Len ( pending ) ]}%
 \@x{\makebox[0pt][r]{\scriptsize 188\hspace{1em}}\@s{27.51} \.{\IN} Deliver (
 latest . at ,\, latest . id ,\, latest . failures )}%
 \@x{\makebox[0pt][r]{\scriptsize 189\hspace{1em}}\@s{16.4} \.{\land} consumed
 \.{'} \.{=} produced \.{\land} pending \.{'} \.{=} {\langle} {\rangle}
 \.{\land} edited \.{'} \.{=} {\FALSE}}%
 \@x{\makebox[0pt][r]{\scriptsize 190\hspace{1em}}\@s{16.4} \.{\land}
 {\UNCHANGED} {\langle} truth ,\, head ,\, based ,\, busy ,\, from ,\, quiet
 ,\, installed ,\, produced ,\,}%
 \@x{\makebox[0pt][r]{\scriptsize 191\hspace{1em}}\@s{89.98} turn ,\,
 toldThisTurn ,\, promptQueued ,\, knows ,\, elapsed ,\, touched ,\, own
 {\rangle}}%
\@pvspace{8.0pt}%
\begin{lcom}{0}%
\begin{cpar}{0}{T}{F}{7.5}{0}{}%
Section: Notification delivery and turn completion
\end{cpar}%
\vshade{5.0}%
\begin{cpar}{1}{F}{F}{0}{0}{}%
A running turn receives at most one notification before another turn starts.
 An idle agent can receive one queued wakeup. Flush publishes eligible
 messages;
 \ensuremath{TurnCompletes} keeps a guarded turn running and permits a clean
 turn to settle.
\end{cpar}%
\end{lcom}%
 \@x{\makebox[0pt][r]{\scriptsize 199\hspace{1em}} Held \.{\defeq} {\IF} turn
 \.{=}\@w{running} \.{\THEN} toldThisTurn \.{\ELSE} promptQueued \.{\lor}
 {\lnot} outbox . waking}%
\@pvspace{8.0pt}%
\@x{\makebox[0pt][r]{\scriptsize 201\hspace{1em}} Flush \.{\defeq}}%
 \@x{\makebox[0pt][r]{\scriptsize 202\hspace{1em}}\@s{16.4} \.{\land} flush
 \.{\land} flush \.{'} \.{=} {\FALSE}}%
 \@x{\makebox[0pt][r]{\scriptsize 203\hspace{1em}}\@s{16.4} \.{\land} {\IF}
 outbox . held \.{\land} {\lnot} Held}%
 \@x{\makebox[0pt][r]{\scriptsize 204\hspace{1em}}\@s{35.71} \.{\THEN}
 \.{\land} outbox \.{'} \.{=} NoMessage \.{\land} knows \.{'} \.{=} outbox .
 failures}%
 \@x{\makebox[0pt][r]{\scriptsize 205\hspace{1em}}\@s{67.02} \.{\land} {\IF}
 turn\@s{0.19} \.{=}\@w{running}}%
 \@x{\makebox[0pt][r]{\scriptsize 206\hspace{1em}}\@s{86.33} \.{\THEN}
 toldThisTurn \.{'}\@s{5.88} \.{=} {\TRUE} \.{\land} {\UNCHANGED}
 promptQueued}%
 \@x{\makebox[0pt][r]{\scriptsize 207\hspace{1em}}\@s{86.33} \.{\ELSE}
 promptQueued \.{'} \.{=} {\TRUE} \.{\land} {\UNCHANGED} toldThisTurn}%
 \@x{\makebox[0pt][r]{\scriptsize 208\hspace{1em}}\@s{35.71} \.{\ELSE}
 {\UNCHANGED} {\langle} outbox ,\, knows ,\, toldThisTurn ,\, promptQueued
 {\rangle}}%
 \@x{\makebox[0pt][r]{\scriptsize 209\hspace{1em}}\@s{16.4} \.{\land}
 {\UNCHANGED} {\langle} truth ,\, head ,\, based ,\, busy ,\, from ,\, quiet
 ,\, installed ,\, produced ,\, consumed ,\,}%
 \@x{\makebox[0pt][r]{\scriptsize 210\hspace{1em}}\@s{89.98} pending ,\,
 report ,\, baseline ,\, wasFailing ,\, told ,\, turn ,\, elapsed ,\, edited
 ,\,}%
 \@x{\makebox[0pt][r]{\scriptsize 211\hspace{1em}}\@s{89.98} touched ,\, own
 {\rangle}}%
\@pvspace{8.0pt}%
\@x{\makebox[0pt][r]{\scriptsize 213\hspace{1em}} TurnStarts \.{\defeq}}%
 \@x{\makebox[0pt][r]{\scriptsize 214\hspace{1em}}\@s{16.4} \.{\land} turn
 \.{=}\@w{idle} \.{\land} turn \.{'} \.{=}\@w{running}}%
 \@x{\makebox[0pt][r]{\scriptsize 215\hspace{1em}}\@s{16.4} \.{\land}
 toldThisTurn \.{'} \.{=} promptQueued \.{\land} promptQueued \.{'} \.{=}
 {\FALSE} \.{\land} flush \.{'} \.{=} {\TRUE}}%
 \@x{\makebox[0pt][r]{\scriptsize 216\hspace{1em}}\@s{16.4} \.{\land}
 {\UNCHANGED} {\langle} truth ,\, head ,\, based ,\, busy ,\, from ,\, quiet
 ,\, installed ,\, produced ,\, consumed ,\,}%
 \@x{\makebox[0pt][r]{\scriptsize 217\hspace{1em}}\@s{89.98} pending ,\,
 report ,\, baseline ,\, wasFailing ,\, outbox ,\, told ,\, knows ,\, elapsed
 ,\, edited ,\,}%
 \@x{\makebox[0pt][r]{\scriptsize 218\hspace{1em}}\@s{89.98} touched ,\, own
 {\rangle}}%
\@pvspace{8.0pt}%
 \@x{\makebox[0pt][r]{\scriptsize 220\hspace{1em}} WokenTurnStarts \.{\defeq}
 promptQueued \.{\land} TurnStarts}%
\@pvspace{8.0pt}%
\@x{\makebox[0pt][r]{\scriptsize 222\hspace{1em}} TurnCompletes \.{\defeq}}%
 \@x{\makebox[0pt][r]{\scriptsize 223\hspace{1em}}\@s{16.4} \.{\land} turn
 \.{=}\@w{running}}%
 \@x{\makebox[0pt][r]{\scriptsize 224\hspace{1em}}\@s{16.4} \.{\land} {\IF}
 Guard}%
 \@x{\makebox[0pt][r]{\scriptsize 225\hspace{1em}}\@s{35.71} \.{\THEN}
 \.{\land} toldThisTurn \.{'} \.{=} {\TRUE} \.{\land} knows \.{'} \.{=}
 report . failures \.{\land} {\UNCHANGED} {\langle} turn ,\, flush {\rangle}}%
 \@x{\makebox[0pt][r]{\scriptsize 226\hspace{1em}}\@s{35.71} \.{\ELSE}
 \.{\land} turn \.{'} \.{=}\@w{idle} \.{\land} flush \.{'} \.{=} {\TRUE}
 \.{\land} {\UNCHANGED} {\langle} toldThisTurn ,\, knows {\rangle}}%
 \@x{\makebox[0pt][r]{\scriptsize 227\hspace{1em}}\@s{16.4} \.{\land}
 {\UNCHANGED} {\langle} truth ,\, head ,\, based ,\, busy ,\, from ,\, quiet
 ,\, installed ,\, produced ,\, consumed ,\,}%
 \@x{\makebox[0pt][r]{\scriptsize 228\hspace{1em}}\@s{89.98} pending ,\,
 report ,\, baseline ,\, wasFailing ,\, outbox ,\, told ,\, promptQueued ,\,
 elapsed ,\, edited ,\,}%
 \@x{\makebox[0pt][r]{\scriptsize 229\hspace{1em}}\@s{89.98} touched ,\, own
 {\rangle}}%
\@pvspace{8.0pt}%
\begin{lcom}{0}%
\begin{cpar}{0}{T}{F}{7.5}{0}{}%
Section: Allowed executions and fairness
\end{cpar}%
\vshade{5.0}%
\begin{cpar}{1}{F}{F}{0}{0}{}%
Next chooses an enabled transition. Spec permits stuttering and requires
 weak fairness for the listed progress actions; this does not impose
 deadlines.
\end{cpar}%
\end{lcom}%
 \@x{\makebox[0pt][r]{\scriptsize 236\hspace{1em}} Next \.{\defeq} Edit
 \.{\lor} OutsideEdit \.{\lor} HeadMoves \.{\lor} QuietElapses \.{\lor} Tick
 \.{\lor} RunEnds \.{\lor} TimesOut}%
 \@x{\makebox[0pt][r]{\scriptsize 237\hspace{1em}}\@s{39.83} \.{\lor}
 ReadResults \.{\lor} Flush \.{\lor} TurnStarts \.{\lor} TurnCompletes}%
\@pvspace{8.0pt}%
\@x{\makebox[0pt][r]{\scriptsize 239\hspace{1em}} Spec\@s{1.46} \.{\defeq}}%
 \@x{\makebox[0pt][r]{\scriptsize 240\hspace{1em}}\@s{16.4} \.{\land} Init
 \.{\land} {\Box} [ Next ]_{ vars}}%
 \@x{\makebox[0pt][r]{\scriptsize 241\hspace{1em}}\@s{16.4} \.{\land} {\WF}_{
 vars} ( QuietElapses ) \.{\land} {\WF}_{ vars} ( Tick ) \.{\land} {\WF}_{
 vars} ( TimesOut )}%
 \@x{\makebox[0pt][r]{\scriptsize 242\hspace{1em}}\@s{16.4} \.{\land} {\WF}_{
 vars} ( ReadResults ) \.{\land} {\WF}_{ vars} ( Flush )}%
 \@x{\makebox[0pt][r]{\scriptsize 243\hspace{1em}}\@s{16.4} \.{\land} {\WF}_{
 vars} ( WokenTurnStarts ) \.{\land} {\WF}_{ vars} ( TurnCompletes )}%
\@pvspace{8.0pt}%
\begin{lcom}{0}%
\begin{cpar}{0}{T}{F}{7.5}{0}{}%
Section: Safety and progress properties
\end{cpar}%
\vshade{5.0}%
\begin{cpar}{1}{F}{F}{0}{0}{}%
Results retain their ordering and checked version. The baseline cannot absorb
 the session\mbox{'}s edits: a non-empty baseline is exactly the failures of
 a version
 at or after the commit that no session edit made or preceded. A failure wakes
 the agent only when its read follows the agent\mbox{'}s own edit, or when
 the check
 timed out, so another process\mbox{'}s changes never wake it. Under
 \ensuremath{Spec} fairness,
 checks end, and each regression, a failure outside the baseline taken,
 reaches
 the agent or waits, without waking it, for its next turn.
 These properties are model-checked within the committed configuration bounds.
\end{cpar}%
\end{lcom}%
 \@x{\makebox[0pt][r]{\scriptsize 256\hspace{1em}} BufferedResultsAreOrdered
 \.{\defeq}}%
 \@x{\makebox[0pt][r]{\scriptsize 257\hspace{1em}}\@s{16.4} \.{\land} \A\, i
 \.{\in} 1 \.{\dotdot} Len ( pending ) \.{:} pending [ i ] . id \.{>}
 consumed \.{\land} pending [ i ] . id \.{\leq} produced}%
 \@x{\makebox[0pt][r]{\scriptsize 258\hspace{1em}}\@s{16.4} \.{\land} \A\, i
 \.{\in} 1 \.{\dotdot} ( Len ( pending ) \.{-} 1 ) \.{:} pending [ i ] . id
 \.{<} pending [ i \.{+} 1 ] . id}%
\@pvspace{8.0pt}%
 \@x{\makebox[0pt][r]{\scriptsize 260\hspace{1em}} ReadsSelectLatestComplete
 \.{\defeq} report . known \.{\implies} report . id \.{=} consumed}%
 \@x{\makebox[0pt][r]{\scriptsize 261\hspace{1em}} ResultsMatchBuilds
 \.{\defeq} report . known \.{\implies} report . failures \.{=} truth [
 report . at ] \.{\lor} report . failures \.{=} \{ TimedOut \}}%
 \@x{\makebox[0pt][r]{\scriptsize 262\hspace{1em}} BaselineNeverAbsorbsEdits
 \.{\defeq}}%
 \@x{\makebox[0pt][r]{\scriptsize 263\hspace{1em}}\@s{16.4} baseline . known
 \.{\implies} baseline . failures \.{=} \{ \} \.{\lor}}%
 \@x{\makebox[0pt][r]{\scriptsize 264\hspace{1em}}\@s{32.8} \.{\land} baseline
 . failures \.{=} truth [ baseline . at ] \.{\land} baseline . at \.{\geq}
 based}%
 \@x{\makebox[0pt][r]{\scriptsize 265\hspace{1em}}\@s{32.8} \.{\land} \A\, w
 \.{\in} own \.{:} w \.{>} baseline . at}%
 \@x{\makebox[0pt][r]{\scriptsize 266\hspace{1em}} RunsEnd \.{\defeq} busy
 \.{\leadsto} {\lnot} busy}%
 \@x{\makebox[0pt][r]{\scriptsize 267\hspace{1em}} EveryEditIsChecked
 \.{\defeq} {\Diamond} {\Box} ( report . known \.{\land} report . at \.{=}
 head )}%
 \@x{\makebox[0pt][r]{\scriptsize 268\hspace{1em}} Regressions \.{\defeq}
 truth [ head ] \.{\,\backslash\,} {\IF} baseline . known \.{\THEN} baseline
 . failures \.{\ELSE} truth [ based ]}%
 \@x{\makebox[0pt][r]{\scriptsize 269\hspace{1em}} WaitsForTurn \.{\defeq}
 turn \.{=}\@w{idle} \.{\land} outbox . held \.{\land} {\lnot} outbox .
 waking \.{\land} Regressions \.{\subseteq} outbox . failures}%
 \@x{\makebox[0pt][r]{\scriptsize 270\hspace{1em}} RegressionsReachAgent
 \.{\defeq} {\Diamond} {\Box} ( Regressions \.{\subseteq} knows \.{\lor}
 TimedOut \.{\in} knows \.{\lor} WaitsForTurn )}%
 \@x{\makebox[0pt][r]{\scriptsize 271\hspace{1em}} OutsideChangesNeverWake
 \.{\defeq}}%
 \@x{\makebox[0pt][r]{\scriptsize 272\hspace{1em}}\@s{16.4} {\Box} [ told
 \.{'} \.{\neq} told \.{\land} told \.{'} . kind \.{=}\@w{failed} \.{\land}
 outbox \.{'} . waking}%
 \@x{\makebox[0pt][r]{\scriptsize 273\hspace{1em}}\@s{30.75} \.{\implies}
 edited \.{\lor} TimedOut \.{\in} outbox \.{'} . failures ]_{ vars}}%
\@x{\makebox[0pt][r]{\scriptsize 274\hspace{1em}}}\bottombar\@xx{}%
\end{document}
